\clearpage
\section{Contemporary solutions: comparison and attribution}
\label{app:contemporary-comparison}

This appendix retains the comparison prepared for version~2.
The later withdrawal of 10.35, which is outside these eleven entries,
is recorded in Appendix~\ref{app:statement-correction}.

\subsection{Scope and sources}

This comparison was prepared by Codex on 19 September 2026 following
Evgeny Khukhro's communication of 18 September. Nine overlaps were
named in the message and two more, 16.9 and 17.101, appeared in its
appended index. Its 17.34 entry is already treated under prior-work
status in Section~\ref{prior:17.34} and is outside the 46 candidates.

We read the eleven papers at the supplied Zenodo DOIs, all attributed
there to Achyuth Jayadevan, and compared them with the arguments in
this paper. For 21.68 we also read Aluna Rizzoli's separate note in the
Notebook repository. The two papers must not be conflated. We retain
the earlier acknowledgment of Zhang and Li for 21.106. Source URLs,
PDF and source-archive hashes, repository metadata, individual review
notes and new checking programs are recorded in
\artifact{paper/reviews/v2-2026-09-19/}.

The assessment below concerns mathematical statements and arguments.
No error was found in the alternatives on this reading, with the
imported theorems and character data retaining their stated roles.
Our new executions cover the specific algebraic identities and finite
character inputs described below. The papers report Lean developments;
we downloaded the eleven accompanying source archives but did not
build them. This comparison is therefore not a kernel verification or
an audit of those projects' full theorem statements and dependencies.
Khukhro reports specialist confirmation of alternative solutions for
21.68, 16.28(a) and 15.89; that report is distinct from our reading.

\subsection{The eleven entries}

\small
\begin{longtable}{@{}p{0.10\textwidth}>{\raggedright\arraybackslash}p{0.80\textwidth}@{}}
\toprule
Problem & Relationship to the argument in this paper\\\midrule
\endfirsthead
\toprule Problem & Relationship to the argument in this paper\\\midrule
\endhead
\bottomrule\endfoot
4.55 & Jayadevan's different example uses \(\Z_{(11)}[\operatorname{PSL}_2(\F_{11})]\),
with total rank 88 and indecomposable rank-44 summands. Our rank is
55,440 over \(\Z_{(5)}[3.A_7]\). Both use completion and rational-character
obstructions to descent~\cite{jayadevan-4-55}.\\[5pt]
14.22 & A different Laurent-polynomial construction gives the radical
obstruction over a two-generated torsion-free polycyclic subgroup of
\(\SL_3(\Z)\). Our coefficient groups are infinitely generated;
the alternative removes that feature~\cite{jayadevan-14-22}.\\[5pt]
15.89 & The same Cayley graph and inverse formula occur, with the stronger
point-spectrum assertion \(\C\setminus\{0\}\) on all functions.
Two-dimensional representations supply the additional eigenfunctions~\cite{jayadevan-15-89}.\\[5pt]
16.9 & The same palindrome-to-reflection reduction, matching recurrence
and cubic operation bound occur. Reconstruction is presented differently.
Qualitative computability and the recurrence also have older sources,
discussed below~\cite{jayadevan-16-9}.\\[5pt]
16.28(a) & The alternative uses a triangular \(\GL_3\) construction
containing the identity. Our trace-one \(\SL_2\) class and central
marker give a different closed set. Both use an algebraically closed
field of characteristic five with a transcendental element~\cite{jayadevan-16-28}.\\[5pt]
17.101 & The HNN-induced-module construction, forest injectivity argument
and cardinal bound agree. Adding regular modules during the iteration
also gives a faithful extension over a nonzero coefficient ring~\cite{jayadevan-17-101}.\\[5pt]
18.76 & The free-factor derivations and commuting-lift obstruction agree.
The characteristic-zero constructions are isomorphic by the explicit
change below, which inverts the kernel. Our argument also permits
positive characteristic~\cite{jayadevan-18-76}.\\[5pt]
21.40 & The bounded-root argument, trace radical and index bound
\((2n+1)^d\), \(d=\dim_{\Q}\operatorname{span}_{\Q}G\), agree.
The alternative separately states the finite-trace-set theorem and a
number-field corollary~\cite{jayadevan-21-40}.\\[5pt]
21.68 & Rizzoli and Jayadevan give the same order-2,592 extension of
Kida's configuration, the same degree-eight representation, and the
same index-two obstruction, in separate accounts~\cite{rizzoli-21-68,jayadevan-21-68}.\\[5pt]
21.106 & Jayadevan gives a different positive formula defining the same
two central generators in the integral Heisenberg group; over a
commutative ring it defines central units. Zhang--Li attribution remains
separate~\cite{jayadevan-21-106,zhang-li-concise}.\\[5pt]
21.121(a) & A different cyclotomic construction gives every prime \(p\),
a locally finite metabelian subgroup of \(\GL_2(\overline{\F_p})\),
and exponent set \((1,\infty)\). Our quaternion/code construction
uses a free product, \(p=2\), and infimum
\(\log24/\log8\)~\cite{jayadevan-21-121a}.\\
\end{longtable}
\normalsize

\subsection{Mathematical checks and prior ingredients}

\paragraph{The smaller projective example, 4.55.}
We reconstructed \(\operatorname{PSL}_2(\F_{11})\) and the displayed
subgroups of orders 12 and 60 from matrices. Induction and removal of
the Steinberg constituent give the three projective characters in the
alternative paper. The fourth follows by subtracting the chosen
completed summand. The ranks are 44, and the involution traces
\(0,-4,4\) distinguish the required modules. GAP's character-table
library agrees with the displayed decomposition data and the irrational
order-five traces of the individual completed projectives. The descent
and indecomposability steps were checked by reading their proofs;
the computation does not replace them.

\paragraph{An explicit comparison of extensions, 18.76.}
Identify the two division closures and write \(g=(x,y)\),
\(h=(x',y')\). Let \(\alpha,\beta\) be the derivations on the first
and second free factors. The cocycles in our coordinates and the
alternative coordinates are respectively
\[
 C_o(g,h)=x\beta(y)\alpha(x'),\qquad
 C_a(g,h)=\alpha(x)y\beta(y').
\]
For \(f(g)=\alpha(x)\beta(y)\), expansion of the derivation laws gives
\[
 (\delta f)(g,h)=g f(h)-f(gh)+f(g)=-C_o(g,h)-C_a(g,h).
\]
Consequently \((r,g)\mapsto(-r+f(g),g)\) is an isomorphism from
the first extension to the second, inducing the identity on the quotient
and negation on the additive kernel. Our exact polynomial check allows
commutation between the two factor algebras only. Thus the difference
between the displayed factor sets does not give a different abstract
group in characteristic zero.

\paragraph{Additional formulas and older algorithms.}
For 15.89 we checked the characteristic polynomial
\(\lambda^2-3s\lambda-4\) and the eigenvector used for every
\(\lambda\ne0\). For 21.106 we checked the determinant identity
that forces the central coordinate to be a unit; the explicit converse
witnesses respect the formula's quantifier order.

For 16.9, Jayadevan's source discussion led us to check
Dahmani--Guirardel\cite[Theorem~4]{dahmani-guirardel-equations} directly.
Their twisted-equation decision procedure gives qualitative computability
of palindromic length, by testing products of at most \(|w|\) palindromes.
The interval recurrence specializes the cancellation-norm algorithm of
Brandenbursky--Gal--K\k{e}dra--Marcinkowski\cite[Section~2.I]{cancellation-norm};
their 2023 erratum establishes word independence in the Coxeter case.
The explicit reflection-length reduction and its application to optimal
palindrome factors are shared with the contemporary paper. Neither
qualitative computability nor the recurrence should be counted as a new
discovery of this experiment.

\subsection{Chronology and what it establishes}

\begin{table}[htbp]
\centering\small
\caption{Artifact chronology in September 2026 (UTC). The first column of
 dates is the first local commit containing the proof file, not a verified
 discovery or upload time. DOI registration is taken from DataCite.}
\label{tab:v2-chronology}
\begin{tabular}{@{}llll@{}}
\toprule
Problem & Local proof-file commit & Zenodo date & DOI registered\\\midrule
4.55 & 11 Sep 08:37 & 14 Sep & 14 Sep 14:08\\
14.22 & 11 Sep 00:48 & 12 Sep & 12 Sep 10:15\\
15.89 & 11 Sep 11:52 & 12 Sep & 12 Sep 10:18\\
16.9 & 11 Sep 12:13 & 12 Sep & 12 Sep 21:08\\
16.28(a) & 11 Sep 00:04 & 13 Sep & 12 Sep 19:39\\
17.101 & 11 Sep 02:22 & 12 Sep & 14 Sep 16:06\\
18.76 & 11 Sep 02:59 & 14 Sep & 14 Sep 14:15\\
21.40 & 11 Sep 07:53 & 14 Sep & 14 Sep 14:43\\
21.68 & 10 Sep 23:00 & 14 Sep & 14 Sep 14:17\\
21.106 & 10 Sep 21:05 & 12 Sep & 12 Sep 21:09\\
21.121(a) & 10 Sep 22:03 & 14 Sep & 14 Sep 14:17\\
\bottomrule
\end{tabular}
\end{table}


The local proof-file commits precede the corresponding DOI registrations.
This comparison concerns those artifacts only. In particular, the Zenodo
publication-date field can differ from record creation and DOI
registration: the 17.101 record lists 12 September, while both technical
timestamps are on 14 September. The 16.28(a) record lists 13 September
but was created and registered on 12 September. Rizzoli's note is dated
11 September; we have not established its first public upload time.
The earlier Zhang--Li artifact dates are discussed in
Section~\ref{subsec:concise}.

These dates do not exclude an earlier draft or discovery, nor do they
establish which sources another AI run could access. The eleven DOI
records expose first-version metadata at retrieval, but we cannot infer
the entire history of the underlying work. The research-repository URL
listed in those records returned HTTP 404 during this comparison, so
we could not inspect its development history.

Our retained structured retrieval index contains no matching URLs for
the eleven supplied DOIs or the compared paper filenames. The exact
patterns and input hash accompany the review. That index omits some
forms of access, including shell downloads, and says nothing conclusive
about training exposure or another run's sources. It is a limited record,
not proof of discovery independence. We accordingly assign no priority
direction and make no copying inference from the similarities.

The historical total remains 46 deadline candidates. The eleven
overlaps are not a uniform category of prior solutions, and the other
35 entries are not thereby established as new. The original mathematical
expositions remain available in version~2, with the individual attribution
notes above; the pre-v2 PDF and source baseline are preserved separately.
